%% kapa: exam=1 round=31 year=2020 period=2 form=A subject=accounting subject_ko=회계학 seq=6 number=46
%% source: Q-Net 감정평가사 기출문제 (cst00309) artlSeq=5207219 file=extracted/1008/5207219/A형/A형/2020년도 제31회 감정평가사 1차 시험_2교시_A형.pdf
%% conversion: kordoc 4.24.51 -> tools/gen1.py (auto); review=proofread-sonnet; flags=none
\begin{question}{46}
(주)감평은 20x1년 초 기계장치(취득원가 ￦1,600,000, 내용연수 4년, 잔존가치 ￦0, 정액법 상각)를 취득하였다. (주)감평은 기계장치에 대해 원가모형을 적용한다. 20x1년 말 동 기계장치에 손상징후가 존재하여 회수가능액을 결정하기 위해 다음과 같은 정보를 수집하였다.
\begin{framedtext}
○ 20x1년 말 현재 기계장치를 처분할 경우, 처분금액은 ￦760,000이며 처분 관련 부대원가는 ￦70,000이 발생할 것으로 추정된다. \par
○ (주)감평이 동 기계장치를 계속하여 사용할 경우, 20x2년 말부터 내용연수 종료시점까지 매년 말 ￦300,000의 순현금유입과, 내용연수 종료시점에 ￦20,000의 기계 철거 관련 지출이 발생할 것으로 예상된다. \par
○ 현재가치 측정에 사용할 할인율은 연 12\%이다. \par
\par{\small\centering\begin{tabular}{|c|c|c|}\hline
기간 & \makecell{단일금액 ￦1의 현재가치\\(할인율=12\%)} & \makecell{정상연금 ￦1의 현재가치\\(할인율=12\%)} \\ \hline
3 & 0.7118 & 2.4018 \\ \hline
\end{tabular}\par}
\end{framedtext}
(주)감평이 20x1년 유형자산(기계장치) 손상차손으로 인식할 금액은? (단, 계산금액은 소수점 첫째자리에서 반올림하며, 단수차이로 인한 오차가 있으면 가장 근사치를 선택한다.)
\begin{choices}
  \item ￦465,194
  \item ￦470,000
  \item ￦479,460
  \item ￦493,696
  \item ￦510,000
\end{choices}
\end{question}
\answer{4}
